% !TEX root = Hogwarts_1347.tex

\chapter{Crécy}

\begin{quote}
\parskip=\baselineskip

\begin{flushright}
\textit{Paris, the 16th day of August, in the year of our Lord 1347}
\end{flushright}

\bigskip

\textit{From Étienne de Beaumont, Conseiller en Chef de la Confrérie des Arts de la Sorcellerie de France, to his esteemed colleague and dear friend Malcolm Dunbar, Chief Councillor to the Wizards' Council of Scotland, Edinburgh.}

\bigskip

Malcolm,

%\bigskip

I trust this letter finds you in good health and that the summer has been pleasant in Edinburgh. I write to you tonight not with the usual matters of commerce and council, but with something that has kept me from sleep for three consecutive nights, and which I fear will keep me from sleep for many more. I ask that you read what follows with the seriousness I know you to possess, and that you forgive the directness of my account. There is no gentle way to present it.

You will have heard of the battle at Crécy. It was fought, in Picardy, between the forces of King Philip and those of the English King Edward. I will spare you the military particulars --- they are the business of Muggle kings and of little consequence to our affairs. What concerns us is what came after.

A soldier was recovered by a French scouting party three days following the engagement. He was the sole survivor found upon the field. He was a common man-at-arms, a man of no particular education or standing, and he was in a state of such severe distress that the physicians who attended him believed he had lost his mind entirely. He had not. I have since spoken with him myself, at considerable length, and I am satisfied that what he witnessed was real.

He told me the following.

He had been wounded during the battle and had lost consciousness. He awoke in darkness, lying among the dead. The field was still. He could not determine how much time had passed, only that it was early into the night. As his eyes adjusted, he became aware of movement among the fallen --- both French and English alike, for the dead do not observe allegiance.

At first he took the shapes for horses, or for men on horseback. They were neither. He described a figure without skin, the raw musculature exposed, the body a grotesque junction of man and horse fused into a single form. The torso of a man rose from the back of a horse's body, but there was no separation between the two --- the flesh and sinew ran continuous, as though the creature had been made wrong from the beginning and had never been corrected. The arms were long, ending in claws of considerable reach. It moved on hooves.

The description is unmistakable. What this soldier encountered is what your countrymen in the Northern Isles call a Nuckelavee.

It was tearing the dead apart.

Not with hunger or purpose that the soldier could discern. The long arms reached down, the claws opened the armour and the bodies beneath with an ease that suggested the resistance of steel and bone was negligible to the creature. It worked methodically, moving from one body to the next. As the soldier watched, paralysed with terror, he counted. There were seven.

Seven Nuckelavees upon the field, moving among the dead of both armies, the hooves striking the helmets and breastplates of the fallen, the claws pulling men apart as one might pull apart wet bread.

A second survivor --- a man the soldier had not known was alive --- made a sound. A breath, or a moan, or the involuntary noise of a body failing. It was enough. The nearest creature turned toward the sound with immediacy. It covered the distance at speed. The soldier heard the impact and the sound that followed. He told me he closed his eyes. When he opened them, the man was in two separate pieces and the creature had already returned to its work.

The soldier understood then that absolute stillness was the only thing preserving his life. He began to move, slowly, using the bodies around him as concealment, dragging himself across the field on his stomach. He told me he crawled for what felt like an hour. He did not look up. He moved toward the sound of running water --- a small stream that crossed the edge of the field.

He reached it. In his exhaustion, or his terror, he lost his footing on the bank and fell into the water. The splash was considerable, and the sound carried.

One of the creatures came. He described it covering the ground between them in seconds. He could hear the hooves on the earth, and then, as the creature reached the margin of the stream, a sound he did not expect --- a violent, agonised shriek. The water thrown up by his fall had struck the creature. Where it touched the exposed flesh, the reaction was immediate and severe. The creature recoiled, writhing.

The soldier, operating now on nothing but the instinct of a man who has found the one thing that hurts the thing trying to kill him, began to throw water. He cupped it in his hands, he kicked it, he hurled it with every motion available to his broken body. The creature retreated. A second one, approaching from behind the first, stopped at the water's edge and would not cross. The soldier kept throwing water until his arms would no longer obey him. Then he collapsed in the stream.

That is where the scouting party found him. Half-submerged, barely breathing, still lying in the running water as though some part of him understood, even in unconsciousness, that it was the only safe place left.

His account was dismissed by the military physicians as battlefield madness. It came to our attention by the usual channels --- the Confrérie maintains observers among the Muggle forces, as you know --- and when the details reached me, I felt a coldness in my chest that has not yet left.

It is widely known among our kind that the appearance of a Nuckelavee is a presage of pestilence. Malcolm, you know the stories as well as I do. Before the great pestilence of Justinian\footnote{The devastating pandemic that struck the Byzantine Empire in 541 AD. It is believed to have killed tens of millions and is one of the deadliest pandemics in recorded history.}, which devastated the Eastern Roman Empire and ravaged the known world from the shores of Persia to Gaul and the Iberian Peninsula, three were seen upon a field of the recently dead — the largest number ever recorded together.

Three, and the known world was brought to its knees. There were seven upon the field at Crécy.

I do not know what is coming, my friend. But I believe something is, and I believe it will be worse than anything the histories record. I am sharing this with you before I present it formally to the Confrérie, because I trust your judgement above that of my own council, and because whatever approaches will not respect the borders between our nations any more than it will respect the borders between our world and the Muggle one.

I urge you to inform the Council and to take whatever precautionary measures you deem appropriate. I will do the same here, once I find the words that will make old men in comfortable chairs listen to the testimony of a common soldier about things they would prefer not to believe in.

Write to me soon. I find I am in need of the counsel of a friend.

\bigskip


With the deepest respect and urgency,


\begin{center}
Étienne de Beaumont
\end{center}

\end{quote}
